\section{SF-E818: positive-dimensional defects of a locally free
semistable
hull}\label{SF-E818-sf-e818-positive-dimensional-defects-of-a-locally-free-semistable-hull}

Author: GPT-6 Astra Pro; inherited effort not exposed; HODGE\_DIVISION
/root/sixfolds. Date: 2026-11-21.

\textbf{Status:} Whole-proof mathematical PASS received from /root and
/root/imprimitive on 2026-11-21 (GPT-6 Astra Pro; inherited effort not
exposed). Root reread the precise SF-E811 and SF-E813 hypotheses.
Imprimitive independently reconstructed every proof step and freshly
read the invoked Fulton/Stacks support-cycle statements. The argument
uses only the already reviewed locally free Hermitian--Einstein result
SF-E811, the factor analysis of SF-E813, and the point-defect theorem
SF-E817. It does not use an admissible Hermitian--Einstein metric on a
singular sheaf.

\subsection{Statement}\label{SF-E818-statement}

Let \(X\) be a polarized complex abelian fourfold with genuine
Rosati-invariant real multiplication by a real quadratic field, of
holomorphic type \((2,2)\). Write its compatible polarization as
\(\theta=\theta_++\theta_-\). Suppose that \(E\) is a torsion-free
coherent sheaf of rank \(r>0\), and that its double dual
\(H=E^{\vee\vee}\) is locally free and admits a filtration by subbundles
with stable locally free quotients of the same \(\theta\)-slope. Suppose
also that

\[
 e^B\operatorname{ch}(E)=r(1-q\theta_+^2/2)(1-q\theta_-^2/2),
 \qquad B\in\operatorname{NS}(X)_{\mathbf Q},\qquad q\in\mathbf Q_{>0}.
 \tag{E818.1}
\]

Then \(E\) is not weakly \(G\)-semiregular for any specified coherent
action of an effective finite translation/tensor subgroup
\(G\subset X\times\widehat X\) for which \(E\) is actually equivariant.

In fact the character equation already forces the quotient \(T=H/E\) to
have dimension zero or to be zero. If it is nonzero, SF-E817 gives the
quantitative kernel bound \(1+3s\), where \(s\) is its number of support
orbits. If it is zero, SF-E813 excludes the locally free sheaf itself.

The locally free stable-factor assumption is imposed on the actual hull
\(H\), not on \(E\). Thus the theorem allows singular stable factors of
\(E\), including those arising by elementary modifications. It makes no
assertion when the reflexive hull is singular, or when a locally free
hull has singular stable Jordan--Hölder factors.

\subsection{1. Strict mixed positivity for every effective
surface}\label{SF-E818-strict-mixed-positivity-for-every-effective-surface}

We first prove a small extension of the effectivity argument in supplied
DEG\_L S-22. For every nonzero effective surface cycle \(Z\) on \(X\),

\[
 \int_X[Z]\theta_+\theta_->0.
 \tag{E818.2}
\]

The forms \(\theta_\pm\) have invariant semipositive representatives
supported on the two complex two-dimensional real-multiplication
eigenspaces \(V_\pm\). Hence their mixed wedge restricts nonnegatively
to every complex surface tangent plane. More explicitly, in unitary
bases it is a positive sum of squared norms of the restrictions of the
mixed holomorphic two-forms in \(V_+^*\otimes V_-^*\). Consequently a
zero integral on an effective cycle forces all these mixed forms to pull
back to zero on a projective resolution of each irreducible component.

At a smooth point of a component, write the Plücker bivector of its
tangent plane as \(u_++u_-\), where \(u_\pm\in\bigwedge^2V_\pm\); the
mixed part vanishes by the preceding conclusion. A decomposable bivector
has square zero, whereas this square is \(2u_+\wedge u_-\). Thus either
\(u_+=0\) or \(u_-=0\), and the tangent plane is exactly \(V_-\) or
\(V_+\). The Gauss map of each irreducible component has finite image
and is therefore constant.

For completeness, constant tangent image in \(V_+\), for example, is
impossible here. Let \(Y\) be a smooth projective resolution of the
component and translate its map to \(X\) to send a base point to zero.
All invariant one-forms annihilating \(V_+\) pull back to zero.
Therefore the homomorphism \(\operatorname{Alb}(Y)\to X\) has tangent
image contained in \(V_+\). Its image is a positive-dimensional abelian
subvariety, since the surface map is nonconstant. But the underlying
real eigenspace of \(V_+\) contains no nonzero rational homology vector:
if \(u\in F\setminus\mathbf Q\) and \(z\ne0\) were such a vector, the
equality \(uz=\sigma_+(u)z\), in rational coordinates, would make the
irrational number \(\sigma_+(u)\) rational. A positive-dimensional
abelian subvariety has nonzero rational homology. This contradiction
proves (E818.2). The proof does not assume that \(X\) is simple.

\subsection{2. The surface part of the defect must
vanish}\label{SF-E818-the-surface-part-of-the-defect-must-vanish}

We use the leading-character/support-cycle identity in the following
exact published form. Fulton, \emph{Intersection Theory}, second edition
(1998), Example18.3.11, printed p.363, identifies the
highest-dimensional component of the Riemann--Roch transformation of a
coherent sheaf with its support cycle weighted by generic lengths.
Theorem18.3(3), printed p.353, identifies that transformation on a
smooth variety with
\(\operatorname{ch}(-)\operatorname{td}(T_X)\cap[X]\). Here \(T_X\) is
trivial. Thus a sheaf \(T\) of support dimension at most \(d\) has
\(\operatorname{ch}_j(T)=0\) for \(j<4-d\), and
\(\operatorname{ch}_{4-d}(T)\) is exactly its effective dimension-\(d\)
support cycle. The generic-length convention is also Stacks Project,
Definition42.10.2, tag \texttt{02QV}. These are published/standard
inputs, not claims to reprove the entire Riemann--Roch theorem.

The canonical double-dual map gives

\[
 0\longrightarrow E\longrightarrow H\longrightarrow T\longrightarrow0,
 \qquad \dim\operatorname{Supp}T\le2.
 \tag{E818.3}
\]

Put \(\alpha=c_1(H)/r=c_1(E)/r\). The first Chern component of (E818.1)
gives \(B=-\alpha\). For any positive-rank sheaf define

\[
 \Delta(F)=\frac{c_1(F)^2}{2\operatorname{rk}F}-\operatorname{ch}_2(F).
\]

Since the quotient in (E818.3) has codimension at least two,

\[
 \Delta(E)=\Delta(H)+\operatorname{ch}_2(T),
 \qquad
 \int_X\Delta(E)\theta_+\theta_-=0.
 \tag{E818.4}
\]

Let \(Q_j\) be the stable locally free quotients in the assumed
filtration of \(H\), and set

\[
 r_j=\operatorname{rk}Q_j,\qquad
 \alpha_j=c_1(Q_j)/r_j,\qquad \delta_j=\alpha_j-\alpha.
\]

The exact Chern-character variance identity is

\[
 \Delta(H)=\sum_j\Delta(Q_j)-\frac12\sum_jr_j\delta_j^2.
 \tag{E818.5}
\]

Each factor has the same slope as \(H\), so
\(\int_X\delta_j\theta^3=0\). The invariant representative of
\(\delta_j\) is therefore primitive for \(\theta\). SF-E811 gives

\[
 \int_X\Delta(Q_j)\theta_+\theta_-\ge0,
 \qquad -\int_X\delta_j^2\theta_+\theta_-\ge0.
 \tag{E818.6}
\]

The second inequality is the elementary mixed sum-of-squares identity
(E811.4), and the first uses a smooth Hermitian--Einstein connection
only on the locally free stable factor \(Q_j\). Consequently the mixed
pairing of \(\Delta(H)\) is nonnegative.

The class \(\operatorname{ch}_2(T)\) is the effective fundamental cycle
of the dimension-two part of \(T\): its coefficients are the positive
generic lengths at the irreducible support surfaces. Terms of smaller
support dimension do not contribute. Equations (E818.2), (E818.4), and
(E818.6) therefore force this surface cycle to be zero. Thus

\[
 \dim\operatorname{Supp}T\le1,
 \qquad
 \int_X\Delta(H)\theta_+\theta_-=0.
 \tag{E818.7}
\]

Every nonnegative summand of (E818.5) has zero mixed pairing. By the
equality case of SF-E811, each \(Q_j\) is projectively flat, with
\(\operatorname{ch}(Q_j)=r_je^{\alpha_j}\). By the equality case of
(E811.4), every \(\delta_j\) is block diagonal for \(V_+\oplus V_-\),
and primitive on each complex two-plane separately. This is precisely
the factor conclusion proved in SF-E813; no splitting of \(H\), and no
equivariance of individual Jordan--Hölder factors, is needed in this
step.

\subsection{3. The curve part of the defect must
vanish}\label{SF-E818-the-curve-part-of-the-defect-must-vanish}

Write \(\delta_j=\delta_{j,+}+\delta_{j,-}\). The preceding primitivity
says

\[
 \delta_{j,+}\theta_+=0,\qquad
 \delta_{j,-}\theta_-=0,
 \qquad \delta_{j,+}^3=\delta_{j,-}^3=0.
 \tag{E818.8}
\]

These are identities of invariant forms on the indicated complex
two-dimensional blocks. Expanding the cube gives

\[
 \delta_j^3=3\delta_{j,+}^2\delta_{j,-}
                +3\delta_{j,+}\delta_{j,-}^2.
\]

Multiplication by \(\theta_++\theta_-\) kills every term, either by
excess degree in one block or by one of the first two identities in
(E818.8). Hence

\[
 \theta\delta_j^3=0,
 \qquad
 \int_X\theta\bigl(e^{-\alpha}\operatorname{ch}(H)\bigr)_3
   =\frac16\sum_jr_j\int_X\theta\delta_j^3=0.
 \tag{E818.9}
\]

The normalized character of \(E\) in (E818.1) has zero degree-three
component. By (E818.7), \(T\) has support dimension at most one, so its
degree-three component is unchanged by multiplication by an exponential.
Additivity in (E818.3) therefore gives

\[
 \operatorname{ch}_3(T)
   =\bigl(e^{-\alpha}\operatorname{ch}(H)\bigr)_3,
 \qquad
 \int_X\theta\operatorname{ch}_3(T)=0.
 \tag{E818.10}
\]

For a sheaf with support dimension at most one,
\(\operatorname{ch}_3(T)\) is its effective fundamental curve cycle with
positive generic-length multiplicities. Every nonzero such curve cycle
has strictly positive degree for the ample polarization \(\theta\). Thus
(E818.10) proves that \(T\) has no one-dimensional component. It is
zero-dimensional or zero, as claimed.

\subsection{4. The remaining cases and exact
scope}\label{SF-E818-the-remaining-cases-and-exact-scope}

If \(T\ne0\), the canonical sequence (E818.3) is precisely the sequence
in SF-E817. The specified coherent finite action on \(E\) canonically
induces coherent actions on \(H\) and \(T\); this functorial statement
is proved there. That theorem gives a nonzero trace kernel, indeed
dimension at least \(1+3s\).

If \(T=0\), then \(E=H\) is a locally free sheaf with the stable locally
free filtration assumed in SF-E813 and the same balanced character.
SF-E813 excludes weak equivariant semiregularity. This proves SF-E818 in
all cases.

The new part relative to SF-E817 is the elimination of all
positive-dimensional defects under an explicit condition on the hull.
There is no assertion that every locally free semistable hull
automatically has locally free stable factors. There is also no use of a
Chern--Weil identity for a singular reflexive sheaf, no extension of a
connection across its singularities, and no reduction of a general
derived object to a sheaf.

The smooth Hermitian--Einstein input has exactly the
rational-polarization scope recorded in SF-E811: Siu, \emph{Lectures on
Hermitian--Einstein metrics for stable bundles and Kähler--Einstein
metrics}, DMV Seminar8 (1987), Chapter1, §1.3--§1.4, pp.6--7;
Proposition1.6, pp.8--9; §6.1--§6.4, pp.35--39. Its proof is not
repeated here. The mixed positivity and its equality case are the
independently checked SF-E811 calculation. The strict surface argument
above fully reproves the needed strengthening of supplied DEG\_L S-22.
The categorical trace input is only that already recorded in SF-E813 and
SF-E817.
